\section{Introduction}\label{sec:introduction}

Two orthonormal bases $\mathcal B$ and $\mathcal C$ of $\C^d$ are
\emph{mutually unbiased} if
\[
 |\langle b,c\rangle|^2=\frac1d
 \qquad(b\in\mathcal B,\ c\in\mathcal C).
\]
Here $\langle u,v\rangle=\sum_j\overline{u_j}v_j$.
For $d\ge2$, let $N(d)$ be the largest number of pairwise mutually unbiased
orthonormal bases in $\C^d$. Our main result is the following.

\begin{theorem}[Computer-assisted]\label{thm:main}
Under the arithmetic and compiler conditions in
Appendix~\ref{sec:rounding}, the fresh complete execution documented in
Section~\ref{execute:observed} certifies that there are no four pairwise mutually unbiased orthonormal bases in
$\C^6$. Consequently, $N(6)=3$.
\end{theorem}

A common unitary transformation makes one basis the coordinate basis.
The second is then the column basis of $H/\sqrt6$, where $H$ is a
complex matrix satisfying
\[
 |H_{ij}|=1,\qquad HH^*=6I.
\]
Such an unscaled matrix is called a complex Hadamard matrix. The
upper-bound problem is to exclude an extension of this pair by
\emph{two} further bases, for every possible $H$. The second basis is
unrestricted: the argument covers arbitrary complex Hadamard matrices,
including those outside previously studied families.

\subsection{Historical context and scope}

Mutually unbiased bases describe projective measurements for which a state
from one basis produces the uniform outcome distribution in every other
basis. This finite-dimensional form of complementarity appears in
Schwinger's unitary-operator framework~\cite{Schwinger1960Unitary}.
Ivanovi\'c developed the state-determination problem and prime-dimensional
constructions~\cite{Ivanovic1981StateDetermination}; Wootters and Fields
proved the statistical optimality of complete sets for state determination
and constructed them in every prime-power dimension
\cite{WoottersFields1989}. One has $N(d)\le d+1$; an operator-space proof
and further prime-power constructions are given in
\cite{BandyopadhyayEtAl2002}. Dimension six is the first dimension outside
the prime-power class.

The tensor-product construction gives three bases in $\C^6$
\cite[Section~4]{KlappeneckerRotteler2004}. Zauner's dimension-six MUB
conjecture asserts that this lower bound is sharp. Its origin is traced
to his 1991 diploma thesis~\cite{Zauner1991Diploma}; see also the historical
account in~\cite[Section~2.3]{McNultyWeigert2026Review}. His doctoral
thesis states the corresponding quantum-design conjecture
\cite[Section~3.3]{Zauner1999}, later translated in~\cite{Zauner2011}.
The equation $N(6)=3$ appears explicitly as Conjecture~1 of
\cite[Section~4]{KlappeneckerRotteler2004}.

Several complementary approaches developed the dimension-six problem.
Grassl proved that the standard--Fourier pair cannot belong to four
mutually unbiased bases~\cite[Theorem~2]{Grassl2004SICMUB6}.
Bengtsson and collaborators studied its
unbiased vectors and the surrounding landscape of Hadamard families
\cite{BengtssonEtAl2007MUBHadamard}. Numerical searches by Butterley and
Hall~\cite{ButterleyHall2007Numerical} and by Brierley and Weigert
\cite{BrierleyWeigert2008} supplied evidence for the conjecture.
Brierley and Weigert also used algebraic equations to enumerate unbiased
vectors for specified Hadamards~\cite{BrierleyWeigert2009}; their
Section~3.4 supplies rigorous approximation bounds for exact input
matrices, while their Section~4.3 explains the limitations introduced by
rounding coefficients in sampled non-affine families.

The closest antecedent of the present proof is the certified
discretization of Jaming, Matolcsi, M\'ora, Sz{\"o}ll{\H o}si and Weiner.
They proved that the standard basis together with any member of the
two-parameter Fourier family cannot be extended to four bases
\cite[Theorem~1.4]{JamingEtAl2009}. Their Section~3 combines phase
intervals, rigorous error bounds, subdivision, and exhaustive
compatibility tests. Jaming, Matolcsi and M\'ora subsequently formulated
an unrestricted discretization program
\cite[Sections~2--4]{JamingMatolcsiMora2010}. The present proof completes
an exhaustive cover of arbitrary second bases. The completion and
displacement estimates described below reduce this cover to a finite
search with certified error bounds.

Recent work on triples retains conjectural sufficient conditions for
$N(6)=3$~\cite[Section~4]{MatolcsiEtAl2026}. The classification preprint
of C\'ardenes Wuttig and Tindall concerns individual order-six Hadamard
matrices and leaves their MUB compatibility unresolved
\cite[Section~VI]{CardenesWuttigTindall2026Classification}. Our proof
constructs a cover of arbitrary second bases directly. Its upper bound
combines that exhaustive cover with local exclusion of two further bases
and the complete execution and arithmetic contract of
Theorem~\ref{execute:finite-theorem}.

\subsection{The covering and exclusion strategy}

The second basis varies throughout the search, so the cover must retain
its possible displacement while testing additional vectors. The proof
first covers every second basis by balls in its real entrywise phases,
then excludes two further bases uniformly throughout each ball.
Figure~\ref{fig:proof-flow} summarizes the five parts of this argument.
\begin{enumerate}
\item Row and column rephasings make the first row and column of $H$
equal to one. The first two rows and columns then contain nine remaining
phases. A symmetry reduction and necessary interval tests cover their
possible values by finitely many seed balls (\cref{sec:seed}).
\item In Stage~A, write $H$ in blocks with row and column sizes $2+4$.
The seed fixes the upper-left, upper-right, and lower-left blocks up to
their phase displacements; the lower-right $4$-by-$4$ block remains to
be completed. At each seed center that passes the frame checks, the
upper-right block and the conjugate transpose of the lower-left block
have rank two, and hence two-dimensional kernels in $\C^4$. Compressing
the unknown block between these kernels gives a $2$-by-$2$ matrix.
The Hadamard equations bound
its distance from $\sqrt6$ times its unitary polar factor. A closed
cover of $U(2)$, the group of $2$-by-$2$ unitary matrices, therefore
accounts for every feasible completion. Uniform Taylor bounds either
exclude a cell or place all its exact completions in certified phase
balls (\cref{sec:completion,sec:phase-cover}). The frames for these kernels
remain fixed at the seed center, which need not itself be feasible.
\item In Stage~B, each ball carries numerical directions approximating
the tangent directions of the Hadamard equations. A verified residual
identity bounds the error in representing every exact solution
displacement by these directions. Covering the tangent coordinates and retaining
this normal bound yields smaller balls covering all the same exact
solutions (\cref{sec:refinement}). No manifold parametrization is required.
\item In Stage~C, a phase-bin search in each final ball covers every
vector unbiased to the first two bases by small vector balls. Make
these vector balls the vertices of two graphs: an $\mathcal O$ edge
retains a possible orthogonal pair, and a $\mathcal U$ edge retains a
possible unbiased pair. Necessary tests remove only impossible edges.
Two additional bases would require two disjoint six-cliques in
$\mathcal O$, with all $36$ cross edges in $\mathcal U$. The finite
search excludes this pattern (\cref{sec:exclusion}).
\item The complete execution accounts for every seed, every refined
ball, and every initially unresolved record. The fallback finishes
with no unresolved record (\cref{sec:execution}).
\end{enumerate}

\begin{figure}[tbp]
\centering
\begin{tikzpicture}[node distance=5mm,
 box/.style={draw=blue!45!black,rounded corners=2pt,fill=blue!3,
 text width=11.7cm,align=center,inner sep=5pt,font=\small},
 arr/.style={-{Latex[length=2mm]},thick,blue!45!black}]
\node[box] (seed) {Arbitrary second basis $H/\sqrt6$\\
Symmetries and closed phase bins $\longrightarrow$ nine-phase seed balls};
\node[box,below=of seed] (polar) {Stage A: all completions of each seed\\
Fixed frames and a polar factor $\longrightarrow$ a complete $U(2)$ cover};
\node[box,below=of polar] (phase) {Stages A and B: certified phase balls\\
Approximate tangent coordinates + bounded normal displacement};
\node[box,below=of phase] (graph) {Stage C: all possible additional vectors\\
Closed vector bins $\longrightarrow$ vector balls $\longrightarrow$ two compatibility graphs};
\node[box,below=of graph] (end) {Every required two-clique pattern excluded\\
Every unresolved record retained, refined, and checked in the fallback};
\draw[arr] (seed)--(polar);\draw[arr] (polar)--(phase);
\draw[arr] (phase)--(graph);\draw[arr] (graph)--(end);
\end{tikzpicture}
\caption{The proof preserves every exact candidate through successive
covers. An exclusion removes a region only after a necessary condition
fails with its certified margin. The execution documented in
Section~\ref{execute:observed} supplies the finite coverage and
terminal-exclusion premises of
Theorem~\ref{execute:finite-theorem}.}
\label{fig:proof-flow}
\end{figure}


The distinction between proposing a numerical object and certifying it
is used throughout. A floating-point center need not be an exact
Hadamard matrix. A computed small eigenspace need not parametrize a
component. A numerical linear solve need not be accurate: its output
is accepted only when a separate bound verifies the required
inequality. Every rejection uses a necessary condition with an
explicit safety margin. A failed certification or an exhausted search
cap retains an unresolved region; it cannot prove nonexistence.

All interval cells may be taken closed. Their overlap, repeated phase
representatives, and redundant balls or graph vertices can weaken the
exclusion but do not omit a possible configuration. The arithmetic
premises are stated, and the finite ranges and rounding margins are
proved, in \cref{sec:rounding}. The accompanying execution and build archives,
identified in \texttt{verification/computation/README.md}, preserve the evidence for
the fresh run documented in Section~\ref{execute:observed}. This proof
uses that run. The separate September 7 execution reported in the source
manuscript has not been authenticated here.

\subsection{The elementary lower bound}

The lower bound uses only the tensor-product construction mentioned
above. We give it explicitly to separate attainment of three bases from
the computer-assisted exclusion of four.

\begin{lemma}[Three mutually unbiased bases in dimension six]
\label{lem:three-mub-six}
There exist three pairwise mutually unbiased orthonormal bases in
$\mathbb C^6$.
\end{lemma}

\begin{proof}
Identify $\mathbb C^6$ with $\mathbb C^2\otimes\mathbb C^3$ in the standard
product coordinates. Define
\[
A_0=I_2,\qquad
A_1=\frac1{\sqrt2}\begin{pmatrix}1&1\\1&-1\end{pmatrix},\qquad
A_2=\frac1{\sqrt2}\begin{pmatrix}1&1\\i&-i\end{pmatrix}.
\]
These matrices are unitary, and for $r\ne s$ every entry of $A_r^*A_s$
has squared modulus $1/2$; the only cross product that is not immediate
from the entries is
\[
A_1^*A_2=\frac12\begin{pmatrix}1+i&1-i\\1-i&1+i\end{pmatrix}.
\]
Let $\omega=e^{2\pi i/3}$, and set
\[
F=\frac1{\sqrt3}(\omega^{jk})_{j,k=0}^{2},\qquad
D=\operatorname{diag}(1,\omega,\omega),\qquad
B_0=I_3,\quad B_1=F,\quad B_2=DF.
\]
The identity $1+\omega+\omega^2=0$ shows that $F$ is unitary; $D$ is
unitary as well. Both $F$ and $DF$ have entries of squared modulus $1/3$.
For their cross product,
\[
(F^*DF)_{jk}=\frac13\sum_{x=0}^{2}\omega^{x^2+(k-j)x}.
\]
If $t=k-j$ is read modulo $3$, the sum on the right is $1+2\omega$ for
$t=0$ and $2+\omega^2$ for $t=1,2$. Each has squared modulus $3$.
Consequently every entry of $B_r^*B_s$ has squared modulus $1/3$ when
$r\ne s$.

For $r=0,1,2$, let the columns of $U_r=A_r\otimes B_r$ form the $r$th
basis. Each $U_r$ is unitary. For $r\ne s$,
\[
U_r^*U_s=(A_r^*A_s)\otimes(B_r^*B_s),
\]
so every entry has squared modulus $(1/2)(1/3)=1/6$. Thus these three
bases are mutually unbiased, and $U_0=I_6$.
\end{proof}

\subsection{A Lie-algebra consequence}

An orthonormal basis $\mathcal B$ of $\C^6$ determines the Cartan
subalgebra $\mathfrak h_{\mathcal B}$ of traceless matrices diagonal
in that basis. It is spanned by the matrices $bb^*-I/6$ for
$b\in\mathcal B$. The Killing form on $\mathfrak{sl}_6(\C)$ is
$K(X,Y)=12\operatorname{tr}(XY)$; two subalgebras are
\emph{Killing-orthogonal} when this form vanishes between them. For
unit vectors $b,c$,
\[
 \operatorname{tr}\bigl((bb^*-I/6)(cc^*-I/6)\bigr)
 =|\langle b,c\rangle|^2-\frac16.
\]
Thus two bases are mutually unbiased exactly when their associated
Cartan subalgebras are Killing-orthogonal. These algebras are closed
under the standard adjoint, and the converse correspondence uses this
same adjoint condition.

\begin{corollary}[Adjoint-closed Cartan subalgebras]
\label{cor:adjoint-cartan}
Under the arithmetic and compiler conditions of
Appendix~\ref{sec:rounding}, using the complete execution documented in
Section~\ref{execute:observed} as in Theorem~\ref{thm:main},
the maximum number of pairwise Killing-orthogonal Cartan subalgebras of
$\mathfrak{sl}_6(\C)$ that are all closed under the same standard adjoint
$X\mapsto X^*=\overline X^{\,T}$ is three.
\end{corollary}

\begin{proof}
By \cite[Theorem~5.2]{BoykinSitharamTiepWocjan2005}, such collections of
Cartan subalgebras correspond to collections of mutually unbiased
orthonormal bases of the same cardinality. Theorem~\ref{thm:main} gives the
upper bound under its stated contract, and Lemma~\ref{lem:three-mub-six}
gives attainment.
\end{proof}
